\documentclass[twocolumn, 12pt]{article}
\usepackage[utf8]{inputenc}
\usepackage{amsmath}
\usepackage{marvosym}
\usepackage{textcomp}
\usepackage{geometry}
\usepackage{array}
\usepackage{wrapfig}
\usepackage{multirow}
\usepackage{graphicx}
\usepackage{float}
\usepackage{wrapfig}
\usepackage{caption}
\usepackage{amssymb}


\geometry{a4paper, margin=.75in}

\begin{document}


\section*{Dynamics}
\subsection*{Equations}


\begin{equation*}
   \textbf{s}_f= \textbf{s}_i + \textbf{v}_it+ \frac{1}{2} \textbf{a}_ct^2   \hspace{15pt} \textbf{v}_f = \textbf{v}_i + \textbf{a}_ct 
\end{equation*}

\begin{equation*}
    \textbf{v}_f^2 = \textbf{v}_i^2+2 \textbf{a}_c (\textbf{s}_f-\textbf{s}_i) \hspace{15pt}      \int_{\textbf{s}_i}^{\textbf{s}_f} \textbf{a}ds=\int_{\textbf{v}_i}^{\textbf{v}_f}  \textbf{v}dv
\end{equation*}

\begin{equation*}
    \int_{\textbf{s}_i}^{\textbf{s}_f}ds =  \int_{t_i}^{t_f} \textbf{v}dt  \hspace{15pt}  \int_{\textbf{v}_i}^{\textbf{v}_f}dv =  \int_{t_i}^{t_f} \textbf{a}dt
\end{equation*}

\begin{equation*}
    \textbf{v}_r = \Dot{r}  \hspace{15pt} \textbf{v}_\theta = r\dot{\theta} \hspace{15pt}\textbf{v}_z = \dot{z}
\end{equation*}

\begin{equation*}
      |\textbf{v}| = \sqrt{\textbf{v}_r^2 + \textbf{v}_\theta^2 +\textbf{v}_z^2} \hspace{15pt} \textbf{a}_z=\Ddot{z}
\end{equation*}

\begin{equation*}
    \textbf{a}_r=\Ddot{r}-r\dot{\theta}^2 \hspace{15pt} \textbf{a}_\theta=r\Ddot{\theta}+2\dot{r}\dot{\theta}
\end{equation*}

\begin{equation*}
    |\textbf{a}|=\sqrt{\textbf{a}_r^2 + \textbf{a}_\theta^2 +\textbf{a}_z^2}
\end{equation*}

\begin{equation*}
    \textbf{a}_n = \frac{\textbf{v}_t^2}{\rho} \hspace{15pt} \textbf{a}_t = \dot{\textbf{v}_t} \hspace{15pt}              \textbf{v}_B = \textbf{v}_A + \textbf{v}_{B/A}
\end{equation*}

\begin{equation*}
    \rho= \frac{(1+(\frac{dy}{dx})^2)^\frac{3}{2}}   {\frac{d^2y}{dx^2}} \hspace{15pt} \Sigma\textbf{F} = m\textbf{a}
\end{equation*}


\begin{equation*}
     U=\textbf{F}s \hspace{15pt} \textbf{F}_s=-k\textbf{s} \hspace{15pt} V_s,U_s = \int_{s_1}^{s_2} \textbf{F}ds
\end{equation*}

\begin{equation*}
    T=\frac{1}{2}m\textbf{v}^2+\frac{1}{2}I_G\omega^2   \hspace{15pt}  T_1+\Sigma U_{2-1}= T_2
\end{equation*}

\begin{equation*}
    T_1+V_1=T_2+V_2 \hspace{15pt} V= mgh + \frac{1}{2}ks^2
\end{equation*}

\begin{equation*}
      V_g=mgh  \hspace{15pt}   1\hspace{3pt} \text{hp} \equiv 550\frac{\text{ft}\cdot \text{lb} }{\text{s}}
\end{equation*}

\begin{equation*}
    P=\frac{dU}{dt}=\textbf{Fv}   \hspace{15pt}  \varepsilon =\frac{\text{power output}}{\text{power input}}
\end{equation*}

\begin{equation*}
    m\textbf{v}_i +\int_{t_1}^{t_2} \Sigma \textbf{F}dt = m\textbf{v}_f
\end{equation*}

\begin{equation*}
    m_a\textbf{v}_{a,i}+m_b\textbf{v}_{b,i}=m_a\textbf{v}_{a,f}+m_b\textbf{v}_{b,f}
\end{equation*}


\begin{equation*}
    e=\frac{\textbf{v}_{b,f}-\textbf{v}_{a,f}}{\textbf{v}_{a,i}-\textbf{v}_{b,i}}  \hspace{15pt}  0 \leq e\leq 1
\end{equation*}

\begin{equation*}
    \textbf{H}=\textbf{r} \times \textbf{p} = I\omega \hspace{15pt} \textbf{r}_i \times \textbf{p}_i = \textbf{r}_f \times \textbf{p}_f
\end{equation*}

\begin{equation*}
    \textbf{r}_i \times \textbf{p}_i + \int \textbf{r}(t) \times \Sigma \textbf{F}(t)dt = \textbf{r}_f \times \textbf{p}_f
\end{equation*}

\begin{equation*}
    \textbf{s}=\theta \textbf{r} \hspace{15pt} \textbf{v}= \Vec{\omega} \times \textbf{r} \hspace{15pt} \textbf{a}_t = \Vec{\alpha} \times \textbf{r}
\end{equation*}

\begin{equation*}
    \int_{\theta_i}^{\theta_f}   \Vec{\alpha}d\theta =\int_{\omega_i}^{\omega_f} \Vec{\omega}d\omega \hspace{15pt} \int_{\theta_i}^{\theta_f} d\theta = \int_{t_i}^{t_f}\vec{\omega} dt
\end{equation*}

\begin{equation*}
    \int_{\omega_i}^{\omega_f}d\omega = \int_{t_i}^{t_f}\Vec{\alpha}dt \hspace{15pt} I=I_G+md^2
\end{equation*}

\begin{equation*}
    \textbf{v}_b = \textbf{v}_a + \Vec{\omega} \times \textbf{r}_{b/a} \hspace{15pt} \textbf{v}_a = \textbf{v}_b +\Vec{\omega} \times \textbf{r}_{a/b} 
\end{equation*}

\begin{equation*}
    \textbf{a}_b = \textbf{a}_a + \Vec{\alpha}\times\textbf{r}_{b/a} -\omega^2\textbf{r}_{b/a} \hspace{15pt}  I_G = mk^2
\end{equation*}

\section*{Undamped Free Vibration}
Standard equation form:
\begin{equation*}
    m \Ddot{x}+kx=0 
\end{equation*}
\noindent
if,
\begin{equation*}
    \omega_n \triangleq \sqrt{\frac{k}{m}}
\end{equation*}

\noindent
then:
\begin{equation*}
    \Ddot{x}+\omega_nx=0
\end{equation*}

\begin{equation*}
    \omega_n=2\pi f \hspace{15pt} f\tau=1
\end{equation*}

\noindent
Simple pendulum period:
\begin{equation*}
    \tau=\frac{2\pi}{\omega_n}=2\pi \sqrt{\frac{l}{g}} 
\end{equation*}

\section*{Confusers}
\begin{center}
\begin{align*}
    H &\equiv \text{angular momentum}\\
    T &\equiv \text{kinetic energy}\\
    U &\equiv \text{work}\\
    V &\equiv  \text{potential energy}\\
    \textbf{v}_{b/a}&\equiv \text{velocity of b, relative to a}\\
    \textbf{r}_{b/a} &\equiv \text{position of b, from a} \\
\end{align*}
\end{center}
\subsection*{Coefficient of Restitution Values}

\noindent
if:
\[
e=
\begin{cases}
1  &\text{the collision is entirely plastic}\\
0 &\text{the collision is entirely elastic}\\
\end{cases}
\]
\begin{center}
    otherwise, the collision is partially elastic
\end{center}




\subsection*{Unit Vectors}


\(\hat{u}_n \equiv \) normal to curvature direction \\
\(\hat{u}_t \equiv\) tangential to motion direction\\ 
\vspace{5pt}\\
\(\hat{u}_\theta \equiv \) transverse direction\\
\(\hat{u}_r \equiv \) radial direction (towards the origin)\\

\subsection*{Diagram Tips}
\underline{Free Body Diagram}
\begin{enumerate}
    \item Apply weight forces
    \item Determine contact points/surfaces
    \item Apply all forces at contact points/surfaces
\end{enumerate}
\underline{Kinetic Diagram}
\begin{enumerate}
    \item Assume motion is in the positive directions
    \item Apply resistance to motion moment(s)
\end{enumerate}

\noindent
\underline{Equations}\\

\noindent
Equate the forces/moments present on the Free Body Diagram to those on the Kinetic diagram, as such:
\begin{equation*}
    \underbrace{\sum\textbf{F}_x}_{\text{Free Body Diagram}} = \underbrace{ \sum \textbf{F}_x}_{\text{Kinetic Diagram}}
\end{equation*}
\begin{equation*}
    \underbrace{\sum\textbf{F}_y}_{\text{Free Body Diagram}} = \underbrace{ \sum \textbf{F}_y}_{\text{Kinetic Diagram}}
\end{equation*}

\begin{equation*}
    \underbrace{\sum\textbf{M}_G}_{\text{Free Body Diagram}} = \underbrace{ \sum \textbf{M}_G}_{\text{Kinetic Diagram}}
\end{equation*}

\vspace{10pt}
\noindent
\underline{Impulse Diagram}
\begin{enumerate}
    \item Draw the initial conditions momentum diagram
    \item Draw the impulse diagram (FBD \(\cdot\) t)
    \item Draw the final conditions momentum diagram
\end{enumerate}
\noindent
Equate the three drawn diagrams according to the following equation:

\begin{equation*}
     \text{Initial Momenta} + \text{Impulse}=\text{Final Momenta}
\end{equation*}




\end{document}